\documentclass[conference]{IEEEtran}

\usepackage{iftex}
\ifPDFTeX
\usepackage[utf8]{inputenc}
\usepackage[T1]{fontenc}
\fi
\usepackage{amsthm}
\usepackage{amsmath,amssymb,bm,mathtools}
\usepackage{newtxtext,newtxmath}
\usepackage{textcomp}
\usepackage{array,booktabs}
\usepackage{graphicx}
\usepackage{siunitx}
\usepackage{cite}
\usepackage[breaklinks=true,hidelinks]{hyperref}
\usepackage{bookmark}
\usepackage[nameinlink,capitalise,noabbrev]{cleveref}
\usepackage[final]{microtype}

\interdisplaylinepenalty=2500
\allowdisplaybreaks
\emergencystretch=2em
\sisetup{detect-weight=true,detect-family=true,per-mode=symbol}

\DeclareRobustCommand{\vct}[1]{\bm{#1}}
\DeclareRobustCommand{\mat}[1]{\bm{#1}}
\newcommand{\R}{\mathbb{R}}
\newcommand{\T}{^{\mathsf T}}
\newcommand{\norm}[1]{\left\lVert#1\right\rVert}
\newcommand{\abs}[1]{\left|#1\right|}
\DeclareMathOperator{\Var}{Var}
\DeclareMathOperator{\Cov}{Cov}
\DeclareMathOperator{\diag}{diag}
\DeclareMathOperator{\clip}{clip}

\newtheorem{theorem}{Theorem}
\newtheorem{proposition}{Proposition}
\newtheorem{corollary}{Corollary}
\theoremstyle{remark}
\newtheorem{remark}{Remark}

\title{Adaptive Integral-State Regularization for Drift-Bounded Inertial Motion Estimation\\
\large Bias--Variance Optimality, Sea-State Scaling, and Horizontal--Vertical Anisotropy (draft)}
\author{%
  \IEEEauthorblockN{Mikhail S. Grushinskiy}
  \IEEEauthorblockA{Independent Researcher, USA, Fair Lawn\\
  Bareboat Necessities GitHub Project\\
  Email: mgrouch@users.github.com}%
}

\begin{document}
\maketitle

\begin{abstract}
Low-frequency drift is the dominant failure mode of acceleration-only marine
motion estimation, yet a hard position-zero constraint removes the long-period
wave motion one is trying to measure.  This paper studies a softer alternative:
augment the inertial chain with an integral state $S=\int p\,dt$ and apply a
virtual $S=0$ measurement whose covariance is adapted to the observed motion.
The virtual measurement is treated explicitly as a \emph{regularizer}, not as a
physical sensor.  First, an Ornstein--Uhlenbeck (OU) similarity argument gives
the natural third-integral scale $\sigma_a\tau^3$, while a normalized Riccati
analysis shows that this scale alone does not determine the correction strength.
A drift-band reduction gives a closed-loop regularization frequency
$\omega_R=(q_{\rm eff}/R_ST_S)^{1/6}$, where $q_{\rm eff}$ is the low-frequency
posterior acceleration-error intensity and $T_S$ is the pseudo-update period.
Because the virtual observation is false whenever real motion has nonzero $S$,
we add the missing signal-distortion term and obtain
$\mathrm{MSE}_p=3q_{\rm eff}/(2\omega_R^3)+4m_{-4}\omega_R^4$.
Its minimizer yields
$r_S^\star\propto q_{\rm eff}^{1/14}m_{-4}^{3/7}T_S^{-1/2}$ and therefore an
axis-specific optimum with no empirical anisotropy constant.  Measurements on
eight stationary calibration seas show that horizontal drift-band acceleration
error is roughly two orders of magnitude above the vertical channel, but the
optimal regularizer changes only modestly: geometric-mean ratios
$r_{S,x}^\star/r_{S,z}^\star=1.138$ and
$r_{S,y}^\star/r_{S,z}^\star=0.861$.  This explains why a historical policy that
made the horizontal constraint much stronger degraded three-dimensional motion
estimation.  Complementary ablations show that the adaptive benefit is carried
primarily by the integral-regularization channel, that the observed amplitude
exponent has a clear minimum near the deployed $r_S\propto\sigma_{aw}$ law, and
that direct OU-prior changes are locally much less important than changes to the
virtual-measurement strength.  The resulting design principle is unusual but
simple: suppress the source of horizontal low-frequency acceleration error; do
not attempt to compensate for it by arbitrarily tightening the integral anchor.
\end{abstract}

\begin{IEEEkeywords}
inertial motion estimation, drift regularization, pseudo-measurement, virtual
measurement, Kalman filter, Ornstein--Uhlenbeck process, adaptive filtering,
heave, anisotropic regularization, bias--variance tradeoff
\end{IEEEkeywords}

\section{Introduction}

Acceleration-only motion estimation contains an unavoidable tension.  Repeated
integration preserves physical displacement in the wave band, but very small
specific-force errors create large low-frequency velocity and position drift.
Marine heave estimators therefore commonly introduce high-pass structure,
model-based restoring dynamics, or virtual measurements in addition to the IMU
\cite{Kuechler2011_HeaveAcceleration,Auestad2013_HeaveStrapdownIMU,Reis2023_DiscreteKalmanHeave}.
For low-cost IMUs the problem is not confined to the vertical axis: attitude
error leaks gravity into horizontal specific force and residual horizontal
accelerometer bias is only weakly observable without an external position or
velocity reference \cite{YurovskyKudinov2025_IMUwaves}.

The estimator studied here uses the augmented linear chain
\begin{equation}
 \dot v=a,\qquad \dot p=v,\qquad \dot S=p,
 \label{eq:chain}
\end{equation}
and periodically applies
\begin{equation}
 y_S=0=S+n_S,\qquad n_S\sim\mathcal N(0,R_S),
 \qquad R_S=r_S^2.
 \label{eq:pseudo}
\end{equation}
The quantity $S$ is the time integral of displacement.  The observation
$S=0$ is not physically true sample by sample; it encodes the long-time belief
that bounded oscillatory motion should not develop an unbounded third integral.
It is therefore better understood as \emph{integral-state regularization} than
as a sensor measurement.

That distinction leads to two separate questions.  First, how should the
regularizer change when the motion time scale and amplitude change?  Second,
should the horizontal axes be regularized more strongly because their inertial
drift is larger?  A dimensional argument answers neither question completely.
The Kalman gain depends on the predicted covariance as well as $R_S$, and a
regularizer strong enough to reject drift also attenuates genuine low-frequency
motion.  The latter effect is absent from a covariance recursion that treats
\eqref{eq:pseudo} as a true observation.

This paper isolates the regularization problem from the surrounding attitude
MEKF and makes five contributions.  It (i) derives the OU similarity scale of
the integral state; (ii) identifies the dimensionless groups of the periodic
Riccati map and the closed-loop drift pole; (iii) restores the signal-distortion
term omitted by the false-measurement covariance model and derives an MSE-optimal
regularizer; (iv) generalizes that optimum per axis and measures its ingredients
on eight directional sea states; and (v) collects the adaptation, sensitivity,
channel-freeze, anisotropy, and transition ablations that test the analytical
predictions.  The implementation vehicle is the repository's OU--III marine
motion estimator, but the reduced results apply to any estimator that anchors a
third integral through a Gaussian virtual measurement.

\section{Integral-State Regularization Model}
\label{sec:model}

For one translational axis, let the latent motion acceleration follow a
stationary OU model
\begin{equation}
 da=-\frac{1}{\tau}a\,dt
 +\sqrt{\frac{2\sigma_a^2}{\tau}}\,dW_t,
 \qquad \Var[a]=\sigma_a^2,
 \label{eq:ou}
\end{equation}
with the chain \eqref{eq:chain}.  OU acceleration is used because it gives a
finite correlation time and an exact discrete transition while remaining a
simple colored-acceleration prior \cite{UhlenbeckOrnstein1930_OU}.  The physical
wave spectrum need not be OU; in the implementation the OU parameters are a
compact local prior fitted to an observed sea-state time scale and acceleration
variance.

The periodic $S=0$ update is applied every $T_S$.  For three-axis motion the
scalar variance becomes
\begin{equation}
 \mat R_S=\diag(r_{S,x}^2,r_{S,y}^2,r_{S,z}^2).
 \label{eq:RS3}
\end{equation}
The important design variable is the standard-deviation scale $r_S$ because the
implementation and the adaptation law are parameterized in those units.

\subsection{The virtual observation is deliberately mis-specified}

If the true displacement is oscillatory but nonzero, then its integral is also
nonzero.  Let $S^{\rm true}$ denote the third integral generated by the desired
motion.  Equation~\eqref{eq:pseudo} is equivalent to a measurement residual
containing the deterministic/stochastic mismatch
\begin{equation}
 d_S=-S^{\rm true}.
 \label{eq:false-input}
\end{equation}
Making $R_S$ smaller increases the authority of the drift anchor, but it also
increases the estimator response to $d_S$.  This is the central bias--variance
trade: one term is unwanted inertial drift, the other is wanted motion that the
virtual measurement contradicts.

\section{Similarity and the Meaning of the Cubic Law}
\label{sec:similarity}

\begin{theorem}[OU similarity scale of the third integral]
\label{thm:similarity}
Under \eqref{eq:ou} and \eqref{eq:chain}, define
\begin{equation}
 \bar t=\frac{t}{\tau},\quad
 \bar a=\frac{a}{\sigma_a},\quad
 \bar v=\frac{v}{\sigma_a\tau},\quad
 \bar p=\frac{p}{\sigma_a\tau^2},\quad
 \bar S=\frac{S}{\sigma_a\tau^3}.
 \label{eq:nondim}
\end{equation}
Then the normalized stochastic dynamics contain neither $\sigma_a$ nor $\tau$.
Consequently the natural coordinate scales of $(a,v,p,S)$ are
\begin{equation}
 \sigma_a,\qquad \sigma_a\tau,\qquad
 \sigma_a\tau^2,\qquad \boxed{\sigma_a\tau^3}.
 \label{eq:natural-scales}
\end{equation}
\end{theorem}

\begin{proof}
Set $t=\tau\bar t$ and use Brownian scaling
$dW_t=\sqrt{\tau}\,d\bar W_{\bar t}$.  Equation~\eqref{eq:ou} becomes
$d\bar a=-\bar a\,d\bar t+\sqrt2\,d\bar W_{\bar t}$.  Successive integrations
then give $d\bar v/d\bar t=\bar a$,
$d\bar p/d\bar t=\bar v$, and $d\bar S/d\bar t=\bar p$.
\end{proof}

A self-similar virtual-measurement standard deviation therefore has the form
\begin{equation}
 r_S=c_R\sigma_a\tau^3.
 \label{eq:cubic}
\end{equation}
Equation~\eqref{eq:cubic} is a coordinate-scale statement.  It does \emph{not}
mean that an unregularized third integral has stationary variance; the OU PSD is
nonzero at zero frequency, so free triple integration diverges.  A finite
third-integral scale requires a low-frequency cutoff, finite horizon, or the
regularizer itself.

\subsection{Periodic Riccati similarity}

Let the IMU interval be $\Delta t$ and let the acceleration observation variance
be $R_a$.  With the diagonal similarity transformation
\begin{equation}
 D=\diag(\sigma_a\tau,\sigma_a\tau^2,
         \sigma_a\tau^3,\sigma_a),
 \label{eq:D}
\end{equation}
the exact discrete periodic Riccati map depends on four dimensionless groups,
\begin{equation}
 \boxed{
 \delta=\frac{\Delta t}{\tau},\qquad
 \rho=\frac{T_S}{\tau},\qquad
 \beta=\frac{R_a}{\sigma_a^2},\qquad
 r=\frac{R_S}{\sigma_a^2\tau^6}.}
 \label{eq:four-groups}
\end{equation}
Thus $R_S\propto\sigma_a^2\tau^6$ fixes only one group.  If $T_S$, $R_a$, or
$\Delta t$ do not scale with the operating point, the normalized Kalman gain is
not invariant even though the units of $R_S$ are correct.

The deployed implementation makes the pseudo-update cadence self-similar away
from safety clamps,
\begin{equation}
 T_S=c_T\tau,
 \label{eq:cadence}
\end{equation}
but renormalizes the per-update variance to preserve the historical continuous
information rate,
\begin{equation}
 r_S=r_S^{\rm base}\sqrt{\frac{T_{S,0}}{T_S}}.
 \label{eq:info-renorm}
\end{equation}
Combining \eqref{eq:cubic}--\eqref{eq:info-renorm} gives the effective input law
\begin{equation}
 r_S\propto\sigma_a\tau^{5/2}
 \label{eq:effective-law}
\end{equation}
whenever \eqref{eq:cadence} is unclipped.  The exponent $5/2$ is therefore a
consequence of cadence scheduling plus information-rate preservation, not a new
natural scale of the third integral.

\section{Drift-Band Riccati Reduction}
\label{sec:riccati}

At frequencies well below the posterior acceleration correlation rate, the
residual acceleration error can be approximated as white with two-sided
intensity $q_{\rm eff}$.  A scalar continuous acceleration estimator gives the
useful reference expression
\begin{equation}
 \mu=\sqrt{\frac{1}{\tau^2}
 +\frac{2\sigma_a^2}{\tau r_a}},
 \qquad
 q_{\rm eff}=2r_a\left(1-\frac{1}{\mu\tau}\right),
 \label{eq:qeff-care}
\end{equation}
where $r_a\simeq R_a\Delta t$ is the high-rate continuous-equivalent
accelerometer noise density.  In the complete inertial estimator,
$q_{\rm eff}$ also contains low-frequency model mismatch, residual bias, and
attitude-to-gravity leakage, so it is ultimately best measured from the
posterior acceleration error rather than inferred from the bench sensor floor.

Reorder the integrated states as $x_d=[S,p,v]^T$ and approximate periodic
measurements by a continuous observation with noise density
\begin{equation}
 r_{S,c}=R_ST_S.
 \label{eq:rsc}
\end{equation}
The continuous algebraic Riccati equation has the closed-form gain
\begin{equation}
 \boxed{
 K_d=\begin{bmatrix}2\omega_R&2\omega_R^2&\omega_R^3\end{bmatrix}^{T},
 \qquad
 \omega_R=\left(\frac{q_{\rm eff}}{R_ST_S}\right)^{1/6}.}
 \label{eq:omegaR}
\end{equation}
The error polynomial is
\begin{equation}
 s^3+2\omega_Rs^2+2\omega_R^2s+\omega_R^3
 =(s+\omega_R)(s^2+\omega_Rs+\omega_R^2).
 \label{eq:poles}
\end{equation}
The virtual-measurement variance therefore sets a genuine low-frequency
regularization bandwidth, albeit through a weak sixth-root dependence.

A fixed normalized pole target $\omega_R\tau=\kappa$ would give
\begin{equation}
 R_S=\frac{q_{\rm eff}\tau^6}{T_S\kappa^6}.
 \label{eq:pole-law}
\end{equation}
This already shows why the natural-coordinate rule \eqref{eq:cubic} is not a
complete tuning law: it contains no $q_{\rm eff}$ and no explicit cadence.
More importantly, even \eqref{eq:pole-law} optimizes only homogeneous drift
rejection.  It still treats $S=0$ as true.

\section{Bias--Variance Optimal Regularization}
\label{sec:mse}

Let the estimator use the gain in \eqref{eq:omegaR} and let the true desired
motion enter through the mismatch \eqref{eq:false-input}.  In the drift-band
reduction the position-error transfer from residual acceleration is
\begin{equation}
 T_{pe}(s)=\frac{s+2\omega_R}{D(s)},
 \qquad
 D(s)=s^3+2\omega_Rs^2+2\omega_R^2s+\omega_R^3,
 \label{eq:Tpe}
\end{equation}
and the transfer from the true $S$ process is
\begin{equation}
 T_{pS}(s)=-\frac{s(2\omega_R^2s+\omega_R^3)}{D(s)}.
 \label{eq:TpS}
\end{equation}
For a white low-frequency acceleration-error intensity $q_{\rm eff}$, the first
term contributes
\begin{equation}
 \mathrm{MSE}_{p,\rm drift}
 =\frac{3q_{\rm eff}}{2\omega_R^3}.
 \label{eq:drift-mse}
\end{equation}
For desired motion concentrated above the regularization corner, define its
negative fourth spectral moment
\begin{equation}
 m_{-4}=\int_{\Omega_{\rm wave}}
 \frac{S_p^{\rm true}(\omega)}{\omega^4}\,d\omega.
 \label{eq:m4}
\end{equation}
The leading low-$\omega_R/\omega$ leakage of \eqref{eq:TpS} then contributes
\begin{equation}
 \mathrm{MSE}_{p,\rm signal}=4m_{-4}\omega_R^4.
 \label{eq:signal-mse}
\end{equation}

\begin{theorem}[MSE-optimal integral regularizer]
\label{thm:mse-opt}
Under \eqref{eq:drift-mse}--\eqref{eq:signal-mse},
\begin{equation}
 \mathrm{MSE}_p(\omega_R)
 =\frac{3q_{\rm eff}}{2\omega_R^3}
 +4m_{-4}\omega_R^4
 \label{eq:mse-total}
\end{equation}
has the unique positive minimizer
\begin{equation}
 \boxed{
 (\omega_R^\star)^7=\frac{9}{32}\frac{q_{\rm eff}}{m_{-4}}.}
 \label{eq:wstar}
\end{equation}
Equivalently, the optimal pseudo-measurement standard deviation is
\begin{equation}
 \boxed{
 r_S^\star
 =\left(\frac{32}{9}\right)^{3/7}
 q_{\rm eff}^{1/14}m_{-4}^{3/7}T_S^{-1/2}.}
 \label{eq:rsstar}
\end{equation}
The minimum position MSE scales as
\begin{equation}
 \mathrm{MSE}_p^\star
 \propto q_{\rm eff}^{4/7}m_{-4}^{3/7}.
 \label{eq:msemin}
\end{equation}
\end{theorem}

\begin{proof}
Differentiating \eqref{eq:mse-total} gives
$-(9/2)q_{\rm eff}\omega_R^{-4}+16m_{-4}\omega_R^3=0$, which is
\eqref{eq:wstar}; the second derivative is positive there.  Substitution of
\eqref{eq:wstar} into
$r_S=\sqrt{q_{\rm eff}/(T_S\omega_R^6)}$, obtained from
\eqref{eq:omegaR}, gives \eqref{eq:rsstar}.  Substitution into
\eqref{eq:mse-total} gives \eqref{eq:msemin}.
\end{proof}

The exponents in \eqref{eq:rsstar} are the important result.  The desired-motion
moment enters with exponent $3/7$, six times the $1/14$ exponent of drift-noise
intensity.  A large increase in low-frequency acceleration error therefore does
\emph{not} imply a proportionally stronger optimal anchor.

\subsection{Amplitude exponent}

For a fixed-shape amplitude sweep, suppose
\begin{equation}
 q_{\rm eff}\propto\sigma_{aw}^{m},
 \qquad m_{-4}\propto\sigma_{aw}^{2}.
 \label{eq:amp-assump}
\end{equation}
Then \eqref{eq:rsstar} predicts
\begin{equation}
 \boxed{
 r_S^\star\propto\sigma_{aw}^{p},
 \qquad p=\frac{m+12}{14}.}
 \label{eq:p-law}
\end{equation}
A sea-independent drift floor ($m=0$) therefore gives $p=6/7\approx0.857$, not
$p=0$.  If drift-band acceleration error grows as
$q_{\rm eff}\propto\sigma_{aw}^2$, then $p=1$ exactly.  The apparently simple
amplitude factor in the deployed law is thus interpretable as a statement about
how posterior low-frequency acceleration error grows with the sea state.

\section{Adaptation-Law Experiments}
\label{sec:adaptation-study}

The repository evaluates the regularizer with paired simulation protocols over
four stationary JONSWAP seas, four PM--Stokes seas, and a controlled transition.
The statistical bundle uses ten predeclared triplets of independent wave-phase,
IMU-noise, and initialization seeds and scores the final \SI{900}{s} of each
\SI{1200}{s} run.  The dedicated sensitivity studies use the same paired-seed
principle.  FixedOracle is a sea-specific, noise-free calibration reference and
is not an online estimator.

\subsection{The amplitude exponent is identifiable}

The cleanest test changes only the exponent $p$ in \eqref{eq:p-law}; all family
members pass through the same nominal $H_s=\SI{1.5}{m}$ operating point.  This
avoids confusing exponent choice with an overall gain change.

\begin{table}[t]
\caption{Amplitude-exponent ablation of the integral regularizer.  OU--III
Adaptive mode, nine scenarios $\times$ ten paired seed triplets.  Vertical
error is percent of $H_s$; 3D is absolute RMS displacement.}
\label{tab:p-ablation}
\centering
\footnotesize
\setlength{\tabcolsep}{3.0pt}
\begin{tabular}{@{}lrrr@{}}
\toprule
Schedule & Z [\%$H_s$] & $\Delta$Z vs $p=1$ & 3D [m] \\
\midrule
$p=0$    & 5.047 & $+0.300\pm0.071$ & 0.712 \\
$p=0.5$  & 4.823 & $+0.076\pm0.025$ & 0.670 \\
$\mathbf{p=1}$ & \textbf{4.747} & --- & \textbf{0.651} \\
$p=1.25$ & 4.779 & $+0.032\pm0.012$ & 0.653 \\
$p=1.5$  & 4.817 & $+0.070\pm0.016$ & 0.660 \\
\bottomrule
\end{tabular}
\end{table}

The minimum is at $p=1$.  Removing amplitude dependence ($p=0$) costs
$0.300$ percentage points of $H_s$, approximately a 6.3\% relative degradation
of the pooled vertical metric, and also worsens 3D RMS.  A controlled amplitude
sweep of the reduced model gives $p=0.854$ for $m=0$ versus $6/7=0.857$,
$p=0.927$ for $m=1$ versus $13/14=0.929$, and $p=1.000$ for $m=2$, confirming
\eqref{eq:p-law}.  Read in reverse, the empirical minimum near $p=1$ is
consistent with a drift-band error intensity growing approximately as
$\sigma_{aw}^2$ over the evaluated envelope.

\subsection{The regularizer channel carries most of the adaptation benefit}

The implementation exposes a $2\times2$ channel-freeze experiment.  The OU
parameters $(\tau,\sigma_{aw})$ and the integral scale $r_S$ are each either
adapted or frozen at the nominal point.

\begin{table}[t]
\caption{OU--III adaptation-channel ablation.  Entries are vertical RMS error
in percent of $H_s$ (mean over ten paired seeds).}
\label{tab:channel-ablation}
\centering
\scriptsize
\setlength{\tabcolsep}{2.3pt}
\begin{tabular}{@{}lrrrr@{}}
\toprule
& Adapt. & $r_S$ only & OU only & Fixed \\
\midrule
$H_s=0.27$ m & 5.13 & 5.13 & 15.82 & 15.72 \\
$H_s=1.50$ m & 4.91 & 4.91 & 4.84 & 4.85 \\
$H_s=4.00$ m & 4.85 & 4.84 & 7.05 & 7.02 \\
$H_s=8.50$ m & 4.03 & 4.02 & 11.28 & 11.14 \\
Transition    & 4.66 & 4.65 & 8.82 & 8.90 \\
\bottomrule
\end{tabular}
\end{table}

Across the off-nominal seas and the transition, adapting $r_S$ with the OU
parameters frozen is essentially indistinguishable from adapting everything.
Conversely, adapting the OU prior while freezing $r_S$ remains close to the
fixed-nominal result.  The direct process-prior adaptation is therefore not the
main source of displacement improvement in this experiment; the adaptive
integral anchor is.

\subsection{One-factor and coupled sensitivity}

A separate local sensitivity sweep perturbs the nominal $H_s=\SI{1.5}{m}$
operating point.  The first three columns in \cref{tab:sensitivity} change one
parameter while the other two are frozen; the final two change $r_S$ together
with the corresponding factor as the deployed law does.

\begin{table}[t]
\caption{Local tuning sensitivity.  Values are vertical RMS error in percent of
$H_s$; standard deviations are omitted here for compactness.}
\label{tab:sensitivity}
\centering
\resizebox{\columnwidth}{!}{%
\begin{tabular}{@{}rrrrrr@{}}
\toprule
Mult. & $\tau$ & $\sigma_{aw}$ & $r_S$ & $\sigma_{aw},r_S$ & $\tau,r_S$ \\
\midrule
0.50 & 4.85 & 4.87 & 5.20 & 5.16 & 9.26 \\
0.75 & 4.85 & 4.85 & 4.87 & 4.87 & 5.45 \\
1.00 & 4.85 & 4.85 & 4.85 & 4.85 & 4.85 \\
1.25 & 4.85 & 4.85 & 4.95 & 4.95 & 5.50 \\
1.50 & 4.85 & 4.85 & 5.12 & 5.10 & 6.89 \\
\bottomrule
\end{tabular}}
\end{table}

With $r_S$ frozen, displacement is almost insensitive to direct changes in
$\tau$ and $\sigma_{aw}$ near the nominal point.  This is expected in a filter
that observes latent acceleration at the IMU rate: the acceleration posterior
is measurement dominated, and $\Delta t/\tau\ll1$.  Once the same parameter
change is propagated into $r_S$, the effect becomes large, particularly for the
cubic coupling to $\tau$.  This is further evidence that adaptation should be
analyzed as a change in regularization bandwidth rather than only as a change in
OU process covariance.

\subsection{Transition and low-motion stress}

\begin{table}[t]
\caption{Selected degradation cases from the paired robustness study.}
\label{tab:stress}
\centering
\footnotesize
\setlength{\tabcolsep}{2.6pt}
\begin{tabular}{@{}llrr@{}}
\toprule
Case & Mode & Z [\%$H_s$] & 3D [m] \\
\midrule
$H_s=0.05$ m & Adaptive & 18.03 & 0.023 \\
360 s ramp & Adaptive & 4.66 & 0.468 \\
360 s ramp & Fixed & 8.88 & 0.720 \\
30 s ramp & Adaptive & 4.87 & 0.486 \\
30 s ramp & Fixed & 9.26 & 0.748 \\
\bottomrule
\end{tabular}
\end{table}

The adaptive law retains most of its performance when the controlled sea-state
transition is shortened from 360 s to 30 s, whereas the nominal fixed point
remains far from the endpoint operating point.  The calm-water case exposes the
opposite limit: when physical motion becomes comparable to the fixed sensor and
attitude-error floor, normalization by $H_s$ necessarily deteriorates.  This is
not repaired by ever-stronger regularization because the low-motion floor is
partly an observability/noise problem rather than a wave-amplitude problem.

\section{Horizontal Versus Vertical Drift}
\label{sec:axes}

The vertical and horizontal channels do not have the same low-frequency error
source.  After attitude compensation, a small tilt error
$\delta\vct\theta$ leaks gravity into horizontal acceleration approximately as
\begin{equation}
 \varepsilon_{a,H}\simeq
 \bigl[g\,\delta\theta_y,\,-g\,\delta\theta_x\bigr]^T
 +\vct b_{a,H}+\vct n_H.
 \label{eq:tilt-leak}
\end{equation}
The vertical channel does not contain the same first-order $g\delta\theta$
term.  Without GNSS or another horizontal kinematic reference, residual
horizontal bias is also substantially less observable.  It is therefore
reasonable to expect $q_{\rm eff,H}\gg q_{\rm eff,z}$.

A common intuitive response is to make $r_{S,x}$ and $r_{S,y}$ much smaller,
thereby making the horizontal $S=0$ constraints stronger.  Theorem
\ref{thm:mse-opt} predicts the opposite caution: $q_{\rm eff}$ enters the
optimal $r_S$ through only the $1/14$ power, while the amount of real motion
being distorted enters through the $3/7$ power.

\subsection{Per-axis optimum}

Applying \eqref{eq:rsstar} independently to axes $i$ and $j$ eliminates all
common constants:
\begin{equation}
 \boxed{
 \frac{r_{S,i}^\star}{r_{S,j}^\star}
 =\left(\frac{q_{{\rm eff},i}}{q_{{\rm eff},j}}\right)^{1/14}
  \left(\frac{m_{-4,i}}{m_{-4,j}}\right)^{3/7}.}
 \label{eq:axis-opt}
\end{equation}
This is an anisotropy law with no new tuning constant.

For each of eight stationary calibration seas, the residual acceleration error
$\hat a_i-a_i^{\rm ref}$ is analyzed below the lower edge of the sea's own
energy band.  A Lorentzian
$2P_i\mu_i/(\mu_i^2+\omega^2)$ is fitted to that drift-band spectrum and its
zero-frequency value is taken as $q_{{\rm eff},i}$.  The desired-motion term is
measured only from the reference displacement spectrum through
\eqref{eq:m4}.  The resulting ratios are in \cref{tab:axis-opt}.

\begin{table}[t]
\caption{Measured per-axis ingredients of the MSE-optimal regularizer.  The
last two columns are the ratios predicted by \eqref{eq:axis-opt}.}
\label{tab:axis-opt}
\centering
\resizebox{\columnwidth}{!}{%
\begin{tabular}{@{}lrrrrrr@{}}
\toprule
& \multicolumn{2}{c}{$q_i/q_z$}
& \multicolumn{2}{c}{$m_{-4,i}/m_{-4,z}$}
& \multicolumn{2}{c}{$r_{S,i}^\star/r_{S,z}^\star$} \\
Sea & x & y & x & y & x & y \\
\midrule
JONSWAP 0.27 & 106.9 & 229.9 & 0.694 & 0.306 & 1.193 & 0.888 \\
JONSWAP 1.50 & 77.2 & 150.1 & 0.729 & 0.272 & 1.191 & 0.819 \\
JONSWAP 4.00 & 32.7 & 51.9 & 0.677 & 0.320 & 1.085 & 0.813 \\
JONSWAP 8.50 & 73.5 & 60.9 & 0.644 & 0.370 & 1.126 & 0.875 \\
PM--Stokes 0.27 & 112.3 & 232.0 & 0.658 & 0.338 & 1.171 & 0.926 \\
PM--Stokes 1.50 & 77.5 & 143.9 & 0.703 & 0.293 & 1.173 & 0.843 \\
PM--Stokes 4.00 & 49.3 & 84.4 & 0.643 & 0.361 & 1.093 & 0.887 \\
PM--Stokes 8.50 & 39.5 & 39.0 & 0.646 & 0.362 & 1.079 & 0.840 \\
\midrule
Geometric mean & 65.4 & 102.3 & 0.674 & 0.326 & \textbf{1.138} & \textbf{0.861} \\
\bottomrule
\end{tabular}}
\end{table}

The horizontal channels carry 65--102 times the vertical drift-band intensity,
yet the optimal regularizer differs by only about $+14\%$ on $x$ and $-14\%$
on $y$.  For example, the factor $65^{1/14}\approx1.35$ from drift noise is
mostly cancelled by $0.674^{3/7}\approx0.84$ from the desired-motion spectrum.
An order-of-magnitude error in $q_{\rm eff}$ moves the optimal $r_S$ by only
about 18\%.  This makes the analytic optimum unusually robust to uncertainty in
the hardest quantity to estimate.

The result also changes where engineering effort should go.  At the optimum,
\eqref{eq:msemin} contains $q_{\rm eff}^{4/7}$, eight times the exponent with
which $q_{\rm eff}$ enters the tuning ratio.  Horizontal RMS is therefore much
more effectively reduced by lowering the actual horizontal drift intensity ---
through better tilt, accelerometer-bias observability, GNSS, or another
horizontal reference --- than by changing $R_S$.

\section{Anisotropy Ablations}
\label{sec:anisotropy}

Earlier versions of the estimator used two horizontal scale factors:
$S_\sigma$ multiplied the horizontal OU acceleration prior, and $\rho_{xy}$
multiplied the horizontal pseudo-measurement standard deviation.  A historical
setting $\rho_{xy}=0.36$ made the horizontal anchor $2.8$ times tighter than the
vertical one.  The hypothesis was that larger horizontal drift required more
regularization.  The experiments below falsify that hypothesis.

\subsection{Two-parameter paired study}

Seven arms separated the OU-prior anisotropy from the pseudo-measurement
anisotropy on four JONSWAP records and five paired seeds.  \cref{tab:aniso-arms}
reports pooled mean percent change relative to the then-deployed point
$(S_\sigma,\rho_{xy})=(1.87,1.00)$; negative is better.

\begin{table}[t]
\caption{Pooled anisotropy ablation relative to $(1.87,1.00)$.}
\label{tab:aniso-arms}
\centering
\footnotesize
\setlength{\tabcolsep}{2.7pt}
\begin{tabular}{@{}lrrrr@{}}
\toprule
$(S_\sigma,\rho_{xy})$ & X [\%] & Y [\%] & Z [\%] & 3D [\%] \\
\midrule
$(1.87,1.87)$ & +2.3 & +27.6 & -0.2 & +13.7 \\
$(1.00,1.00)$ & -4.2 & -1.1 & +0.1 & -2.6 \\
$(1.00,1.87)$ & -12.8 & +26.8 & -0.1 & +6.4 \\
$(0.80,1.00)$ & -0.9 & -1.7 & +1.1 & -1.4 \\
$(0.80,0.80)$ & +5.3 & -7.1 & +1.3 & -0.2 \\
$(1.87,0.53)$ & +21.6 & -8.8 & +1.1 & +8.7 \\
\bottomrule
\end{tabular}
\end{table}

Changing $\rho_{xy}$ mainly trades $x$ error against $y$ error.  Loosening the
horizontal anchor helps one axis and hurts the other because the directional
records place different fractions of physical wave motion on world $x$ and
world $y$.  Three-dimensional RMS loses in either direction.  A single
``horizontal'' multiplier is therefore too coarse to solve an axis-dependent
signal-versus-drift problem.

\subsection{Horizontal OU-prior scale sweep}

After the two-dimensional study isolated most of the mis-specification on the
OU-prior side, $S_\sigma$ was swept with the regularizer kept isotropic over all
eight stationary records and five seeds.  \cref{tab:ssigma} reports pooled
percent change relative to the old $S_\sigma=1.87$ point.

\begin{table}[t]
\caption{Horizontal OU-prior scale sweep with isotropic integral regularization.}
\label{tab:ssigma}
\centering
\footnotesize
\setlength{\tabcolsep}{3.0pt}
\begin{tabular}{@{}rrrrr@{}}
\toprule
$S_\sigma$ & X [\%] & Y [\%] & Z [\%] & 3D [\%] \\
\midrule
0.60 & +9.59 & -1.42 & +3.45 & +4.62 \\
0.80 & -0.42 & -0.88 & +0.79 & -0.64 \\
0.90 & -1.75 & -0.69 & +0.37 & -1.19 \\
\textbf{1.00} & \textbf{-2.11} & \textbf{-0.54} & +0.16 & \textbf{-1.28} \\
1.10 & -2.04 & -0.42 & +0.05 & -1.19 \\
1.20 & -1.79 & -0.32 & -0.00 & -1.02 \\
1.50 & -0.87 & -0.13 & -0.03 & -0.49 \\
\bottomrule
\end{tabular}
\end{table}

The minimum is broad and centered at one, supporting the current isotropic OU
acceleration prior.  Together with \cref{tab:axis-opt}, the result is more
interesting than simply ``anisotropy did not help'': the horizontal disturbance
is dramatically anisotropic, but the MSE-optimal \emph{regularizer} is nearly
isotropic because the desired horizontal motion is anisotropic as well.

\section{Additional Ablations and Mechanistic Checks}
\label{sec:checks}

Several secondary experiments help separate the regularization mechanism from
implementation artifacts.

\subsection{Covariance re-alignment is not the adaptation gain}

The adaptive implementation historically re-aligned the posterior $a_w$
marginal to the stationary OU covariance.  Matched
Adaptive/AdaptiveHeldCovariance runs switch that periodic operation off while
keeping the same tuning trajectory.  Across the main scenarios the vertical
metric changes only at the hundredth-of-a-percentage-point level.  The useful
adaptive effect therefore does not come from periodically inflating
$P_{a_wa_w}$; it comes from the operating point, chiefly $r_S$.

\subsection{The wave-band period must be estimated in the correct band}

A former tuner derived $\tau$ from an acceleration-band frequency.  Because
acceleration weights elevation spectra by $\omega^4$, that frequency barely
changed as the physical wave period grew.  A replacement estimator recovers the
wave-band zero-crossing period from a levelled vertical acceleration signal via
high-pass filtering and paired leaky integrators.  On four JONSWAP records the
estimated versus reference $T_z$ values were approximately
$2.60/2.51$, $4.37/4.48$, $7.09/6.60$, and $8.43/8.55$ s.  Feeding the same
period estimator raw body-$Z$ acceleration instead of a measurement-only
levelled signal degraded OU--III vertical RMS by a geometric factor of 2.51 and
3D RMS by 2.24 across the eight records.  A private complementary attitude
observer reproduced the MEKF-levelled result to within 0.2\%, opening the tuner
loop without sacrificing the wave-band estimate.

This matters because $\tau$ enters the regularizer with a high power.  A
frequency estimator that is excellent for phase tracking can still be a poor
sea-state time-scale estimator, and cubing or otherwise strongly weighting that
error produces a large regularization error.

\subsection{The bias--variance trade is visible directly}

On a representative $H_s=\SI{1.5}{m}$ deterministic replay, changing only the
integral-regularizer coefficient away from its nominal value increased 3D RMS
from about $0.273$ m to $0.354$ m for a much tighter anchor and to $0.391$ m for
a much looser one.  The same experiment showed why restoring the historical
horizontal factor $\rho_{xy}=0.36$ was not a solution: it reduced one horizontal
axis slightly while increasing the other enough to raise 3D RMS.  With sensor
noise disabled, the softer integral-anchor architecture becomes comparatively
better, confirming that the horizontal penalty is fed by real low-frequency
inertial error rather than by unavoidable attenuation of the wave itself.

\section{Design Implications}
\label{sec:design}

The combined analysis suggests a hierarchy for adaptive drift regularization.

First, separate three concepts that are easily conflated: the natural coordinate
scale $\sigma_a\tau^3$, the Kalman/observer bandwidth $\omega_R$, and the
MSE-optimal trade between drift suppression and physical-signal distortion.
Only the last one answers how strong the virtual measurement should be.

Second, schedule the pseudo-update cadence explicitly.  A change in update
frequency changes the continuous information rate even when $R_S$ is held
fixed.  Any adaptation law stated only as a power of $\tau$ is incomplete unless
$T_S(\tau)$ is also stated.

Third, estimate the time scale from the physical wave band rather than from the
high-frequency-weighted acceleration spectrum.  For a third-integral
regularizer, time-scale error is amplified strongly by the tuning law.

Fourth, do not derive anisotropy from process-noise amplitude alone.  The
correct per-axis decision uses both $q_{{\rm eff},i}$ and $m_{-4,i}$ through
\eqref{eq:axis-opt}.  In the present data, enormous horizontal drift does not
justify an enormous horizontal correction.

Finally, if horizontal RMS remains excessive while \eqref{eq:axis-opt} places
$r_{S,H}$ near the isotropic value, the limiting problem is observability.
Better attitude, accelerometer-bias estimation, GNSS velocity/position,
water-speed constraints, or other horizontal references attack
$q_{\rm eff,H}$ directly.  Tightening $R_S$ merely moves along the
bias--variance curve and eventually destroys desired motion.

\section{Scope and Limitations}
\label{sec:limitations}

The analytical reduction assumes a scalar linear inertial chain, a separated
low-frequency drift band, and a pseudo-update cadence fast enough to admit the
continuous-information approximation.  The full marine MEKF contains attitude,
bias, and cross-covariance couplings that are not represented in the closed-form
CARE.  The role of those couplings is partly absorbed into the measured
$q_{\rm eff}$, but the theorem should still be read as a reduced design result,
not an exact solution of the complete nonlinear estimator.

The per-axis $q_{\rm eff}$ and $m_{-4}$ measurements come from synthetic
stationary calibration seas with known reference acceleration and displacement.
A real deployment does not have those references online.  Their present value is
to establish the direction and scale of the optimum and to show that a single
large horizontal tightening is unjustified.  An online anisotropic law would
need reference-free estimators or auxiliary navigation measurements and should
be validated separately.

The ten-seed validation bundle and the dedicated anisotropy studies were run at
different stages of the implementation history.  In particular, the committed
full validation bundle was generated before the final horizontal OU-prior factor
was changed from 1.87 to the currently deployed isotropic value 1.0.  The
adaptation-channel and robustness tables therefore establish the behavior of
the tested regularization mechanisms, while the dedicated five-seed anisotropy
sweep is the direct evidence for the current $S_\sigma=1$ choice.  A future full
regeneration should collapse that historical distinction.

No claim is made that the $S=0$ virtual observation is optimal for all marine
motions.  Crossing seas, sustained maneuver acceleration, long-period swell,
and non-wave translation can change $m_{-4}$ or violate the zero-mean integral
model.  The bias--variance formulation is useful precisely because such motion
appears as an explicit signal-distortion term rather than being hidden inside a
Kalman covariance.

\section{Reproducibility}

The results summarized here are already represented by executable studies in
the repository.  The full paired adaptation study is produced by
\texttt{tools/ou\_validation.py}; local sensitivity and transition stress are
produced by \texttt{tools/ou\_robustness.py}; the amplitude-exponent family by
\texttt{tools/ou\_rs\_law\_ablation.py}; the bias--variance calculation by
\texttt{tools/rs\_law\_bias\_variance.py}; and the per-axis quantities in
\eqref{eq:axis-opt} by \texttt{tools/ou\_axis\_rs\_optimum.py}.  The committed
statistical bundles retain seed triplets, source/input hashes, raw rows, derived
summary tables, and paired effects.  The anisotropy experiment is documented in
\texttt{docs/ou-iii-anisotropy-consistency.md}.

\section{Conclusion}

Integral-state pseudo-measurements provide a useful middle ground between free
inertial integration and hard position/velocity zeroing.  Their tuning cannot be
reduced to dimensional consistency.  The natural third-integral scale
$\sigma_a\tau^3$ explains the state units; the periodic Riccati map explains the
closed-loop regularization bandwidth; and the fact that $S=0$ is a false
measurement introduces the missing physical-signal distortion term.  Combining
the last two yields a compact optimum,
\[
 r_S^\star\propto q_{\rm eff}^{1/14}m_{-4}^{3/7}T_S^{-1/2},
\]
which explains both the measured amplitude law and the unexpectedly weak need
for horizontal-versus-vertical regularizer anisotropy.

The most striking experimental result is that horizontal drift intensity can be
two orders of magnitude above the vertical channel while the optimal
regularizer remains close to isotropic.  The reason is not that horizontal drift
is unimportant; it is that regularization strength is a poor lever against a
disturbance that also competes with real horizontal wave motion.  At the optimum,
residual error is far more sensitive to the disturbance itself than the tuning
ratio is.  For drift-bounded inertial motion estimation, that shifts the design
priority from ``regularize the noisy axis harder'' to ``make the low-frequency
acceleration error observable and smaller.''

\bibliographystyle{IEEEtran}
\bibliography{w3d}

\end{document}
